% !PRINTED 206
\weberpage{206}
\begin{center}
\S.\,67.

\gesperrt{The full system of invariants of the binary\\
biquadratic form.}
\end{center}

We wish further to prove that with the constructs $A$, $B$ the system of
independent invariants of the binary biquadratic form is exhausted.
But whereas one usually contents oneself with the proof that every
invariant is an integral rational function of $A$, $B$, we wish, as we
have already done in \S.\,63 for the covariants of the cubic form, also
to touch on the question of the numerical coefficients, and here a new
phenomenon shows itself.

The invariant $D$, for example, is indeed by \S.\,64, (11) expressed rationally
in terms of $A$, $B$. But the numerical coefficients are not integers; rather
they have the denominator 27, although the coefficients in the expanded function
$D$ are all integers.

We therefore now consider as integral invariants integral rational homogeneous
functions of degree $\mu$ of the five variables $a_0, a_1, a_2, a_3, a_4$
\weq{1}{
I(a_0, a_1, a_2, a_3, a_4) = I(a),
}
whose coefficients are integers, and to which the invariant property
\weq{2}{
I' = r^{2\mu} I
}
belongs, when $I'$ or $I(a')$ is the same function of the coefficients of a
transformed function.

Such invariants are $A$, $B$, $D$, between which there holds the relation
\weq{3}{
27 D = 4A^3 - B^2,
}
and our goal is to prove that all integral invariants are integral rational
functions, with integer coefficients, of these three.

If we form the function $I'$ in the normal form (4) of \S.\,65, then
we first obtain from (2), since $r = 1$,
\[
I(a) = I'(0, a_1', a_2', a_3', 0),
\]
and this function $I$ can depend only on $a_2'$ and the product $a_1'a_3'$;
for the substitution
%CONT>>
